% !TeX root = ../MAgent_Theory_2.0_companion.tex
\chapter{Coupling Statistical Change Across Contexts}
\label{comp:contextual-coupling}

A context-local law flow can change a statistical section without comparing its values at neighboring contexts. Giving the connection a dynamical role requires an energy that reads those comparisons. The construction here penalizes the change of collective statistical shape after a common frame motion has been removed. Its gradient then adds a neighboring-context force to the reciprocal law dynamics. On a finite contextual graph this force is a sum of aligned geodesic displacements. On a continuum base it becomes a diffusion of shape coupled to local frame and routing equations.

\section{The statistical material and its local forces}

At each context, agent $i$ has normalized laws $q_i$ and $s_i$ on standard Borel state and model spaces. Restrict its joint law to $\zeta_i=q_i\otimes s_i$, with each factor in a finite-dimensional smooth statistical family. Assume quadratic-mean differentiability, centered square-integrable directional scores $\ell_U$, and a smooth positive-definite Fisher tensor after redundant coordinates have been removed. The family need not be Gaussian. Fixed bimeasurable comparisons place the agent laws in compatible target spaces, giving $p=(p_i)_i\in\mathcal N$. Use the actual normalized product experiment and its metric:
\begin{equation}
 P_p=\bigotimes_i p_i,\qquad
 \mathsf G_p(U,W)=\sum_i g^F_{p_i}(U_i,W_i),\qquad
 g^F_p(U,W)=\E_p[\ell_U\ell_W].
 \label{comp:ctx:metric}
\end{equation}
Product scores add and their cross expectations vanish. Within each agent, the declared product restriction gives the sum of belief and model Fisher metrics. An unrestricted joint law would also carry dependence directions. The restriction is part of this dynamical model; the unrestricted law flow need not preserve it.

Choose proper references $r_i(dm)$, normalized evaluators $G_i(m,dk)$, and normalized channels $\mathcal K^q_{ij}[p]$, $\mathcal K^s_{ij}[p]$. The channel rules may depend on all current laws but not on source rows. Put $R_i=r_iG_i$, $D^q_{ij}=\KL(q_i\Vert q_j\mathcal K^q_{ij}[p])$, and $D^s_{ij}=\KL(s_i\Vert s_j\mathcal K^s_{ij}[p])$. For finite nonempty source sets, positive normalized prior rows $\pi_i^q,\pi_i^s$, and temperatures $T_q,T_s>0$, the local potential is
\begin{equation}
 \begin{aligned}
 \mathcal E(p,\beta,\gamma)={}&\sum_i\KL(q_i\otimes s_i\Vert R_i)\\
 &+\sum_i\left[\sum_j\beta_{ij}D^q_{ij}
                         +T_q\KL(\beta_i\Vert\pi_i^q)\right]\\
 &+\sum_i\left[\sum_j\gamma_{ij}D^s_{ij}
                         +T_s\KL(\gamma_i\Vert\pi_i^s)\right].
 \end{aligned}
 \label{comp:ctx:potential}
\end{equation}
All displayed terms must be finite and sufficiently smooth on the chosen chart. Smooth Fisher geometry alone does not imply smooth KL costs. Data are fixed in evolution time; every declared law dependence is differentiated. In a justified density calculation this includes
\begin{equation}
 d\KL(a\Vert b)=\int\log(a/b)da-\int(a/b)db,\qquad
 d(q_j\mathcal K[p])=(dq_j)\mathcal K[p]+q_jD\mathcal K[p](dp).
 \label{comp:ctx:roles}
\end{equation}
Thus a source moves in response to its receivers as well as its own costs. This is the reciprocal law potential. The product experiment in \cref{comp:ctx:metric} specifies its mobility; it does not make this potential the variational free energy of a fixed population model. Own-block per-agent updates belong to the per-agent objective, omit some of \cref{comp:ctx:roles}, and have no general total-energy dissipation identity. They also need their own regularity hypothesis: smoothness of the sum of all costs does not make each agent's own-block gradient locally Lipschitz, because law-dependent channels can cancel nonsmooth parts between receivers.

One finite-dimensional Lie group $G$ acts diagonally on $\mathcal N$ by sample pushforward, with its respective actions on the two sectors. Assume the family is stable under this action. Sample bijections preserve Fisher inner products, so the action is isometric. Write $J_p\xi=\left.\partial_a\right|_{a=0}\exp(a\xi)\cdot p$ for its fundamental tangent. Comparisons remain fixed in the evolution. Making them state dependent would change both the metric pullback and the variation.

\section{A finite graph of contextual comparisons}

Let the contextual vertices have masses $m_v>0$ and the unoriented edges have weights $w_e>0$. Agent $i$ at one context is compared with agent $i$ at the next; this identification of agents across contexts is part of the model. Choose one orientation $e:u\to v$ and an auxiliary group element $U_e$ carrying target coordinates from $u$ to $v$. Define
\begin{equation}
 c_e(p_u,p_v,U_e)=\frac12d_{\mathcal N}(p_v,U_e\cdot p_u)^2,
 \qquad \mathcal C=\sum_e w_e c_e.
 \label{comp:ctx:edge}
\end{equation}
The distance is intrinsic to the admitted Fisher family. Reversing this geometric edge uses $U_e^{-1}$ and leaves its cost unchanged. Independent attending links between agents at one context remain separate data; their opposite directions need not be inverses. Finite edge transports are group elements rather than connection one-forms; they transform like the link variables of lattice gauge theory \citep{Wilson1974}, and discrete connections have their own geometric theory \citep{LeokMarsdenWeinstein2005}.

\hypothesisheading{A regular finite alignment branch}{comp:ctx:finite-branch}
For every edge, choose a stationary transport where the compared endpoints have a unique minimizing geodesic, squared distance is smooth, and the transport Hessian is positive definite in group coordinates. Continue only the resulting local minimizing branch $U_e^*(p_u,p_v)$ while these conditions hold. The continuation lives on the set of regular critical pairs $(p,U)$, not on law space: carried around a closed loop of laws, a branch can return to a different transport, and it can stop being the global minimizer without violating any condition.

The first variation of half squared distance at an endpoint is minus the geodesic logarithm toward the other endpoint. Varying $U_e$ as $\exp(a\xi)U_e$ therefore gives
\begin{equation}
 \mathsf G_{U_ep_u}\bigl(\operatorname{Log}_{U_ep_u}p_v,
                                      J_{U_ep_u}\xi\bigr)=0.
 \label{comp:ctx:finite-normal}
\end{equation}
The residual is orthogonal to a common frame motion. The implicit-function theorem supplies the smooth branch in \cref{comp:ctx:finite-branch}. Its transport derivative vanishes, so differentiation of $\mathcal C_*(p)=\mathcal C(p,U^*(p))$ keeps only endpoint derivatives. A local minimizing branch need not realize the global infimum over $G$.

Invariance of the edge cost under a common motion of both endpoints, combined with stationarity in $U_e$, makes the endpoint force at $p_v$ Fisher orthogonal to the group orbit: the connecting geodesic is horizontal. The contextual force therefore never rotates frames. Where the selected branch is the global minimizer and the action is free and proper, half the minimized squared distance is half the squared distance between orbits, so $\mathcal C_*$ is a graph Dirichlet energy of the shape map into the quotient $\mathcal N/G$. For planar location laws this is the orthogonal Procrustes problem \citep{Schonemann1966} and the Procrustes distance of shape theory \citep{Kendall1984}. For full Gaussian laws no closed-form Fisher--Rao distance is known \citep{PineleStrapassonCosta2020}; a declared invariant divergence such as relative entropy then gives a computable edge cost with the same local proofs and the same infinitesimal normal equation, at the price of a changed model.

Give each interior source row its categorical Fisher metric and use
\begin{equation}
 \begin{aligned}
 \mathbb G&=\sum_v m_v\left(\mathsf G_{p_v}
                      +\sum_i g^{\mathrm{row}}_{\beta_{iv}}+\sum_i g^{\mathrm{row}}_{\gamma_{iv}}\right),\\
 \mathcal F_*&=\sum_v m_v\mathcal E_v+\lambda\mathcal C_*,\qquad\lambda\geq0.
 \end{aligned}
 \label{comp:ctx:finite-total}
\end{equation}
Up to the constant $\sum_vm_v$, this metric is the Fisher metric of an experiment that first samples the context with probability proportional to $m_v$, then samples its product laws and independent source labels. Its natural-gradient velocity \citep{Amari1998} is
\begin{equation}
 \dot p_v=-\operatorname{grad}_{\mathsf G}\mathcal E_v
 +\frac{\lambda}{m_v}\sum_{e=\{u,v\}}w_e
       \operatorname{Log}_{p_v}(U^*_{v\leftarrow u}\cdot p_u).
 \label{comp:ctx:finite-flow}
\end{equation}
Here the reversed orientation uses the inverse transport. The contextual force points toward aligned neighbors. The masses cancel from the local force, so zero coupling recovers the chosen unit-mobility restricted-family flow of the reciprocal potential.

For either sector, write its row as $\eta$, its temperature as $T$, and $h_j=D_j+T\log(\eta_j/\pi_j)$. Its equation and the total dissipation are
\begin{equation}
 \dot\eta_j=-\eta_j\left(h_j-\sum_l\eta_lh_l\right),\qquad
 \frac{d\mathcal F_*}{dt}=-\|\dot p,\dot\beta,\dot\gamma\|_{\mathbb G}^2.
 \label{comp:ctx:finite-dissipation}
\end{equation}
Pairing the centered row force with a mass-zero row variation gives minus its energy derivative. Pairing the full velocity with $d\mathcal F_*$ proves the second identity. With $C^2$ local potentials, the reduced field is $C^1$ in ordinary law and interior-row coordinates, so the integral equation is a contraction on a sufficiently short interval. This gives a unique local solution; carried together with its transports, it has a maximal continuation on the regular critical set. Energy decrease does not prevent escape from that set, and the reduced energy is the minimized cost only while the branch stays globally minimizing.

Free action alone does not supply the branch hypothesis. For two planar Gaussian location means, let $SO(2)$ act on both columns and take $X_\epsilon=(e_1,(1+\epsilon)e_2)$ and $Y=(e_1,-e_2)$. Both configurations have trivial stabilizer near zero, but direct rotation gives
\begin{equation}
 \min_R\frac12\|X_\epsilon-RY\|_F^2
       =2+\epsilon+\frac12\epsilon^2-|\epsilon|.
 \label{comp:ctx:cusp}
\end{equation}
The minimizer switches from $I$ to $-I$, with every rotation minimizing at zero. The reduced energy is not differentiable there, and its descent directions depend on which minimizing rotation is used, so a gradient flow of the minimized energy is not unique from that point. Conversely, for a free action at $p_0$ that is compact, or proper on a complete family, agreement of the contextual endpoints at $(p_0,p_0)$ has a unique minimizer and positive Hessian $\mathsf G(J\cdot,J\cdot)$; globally minimizing regular alignments exist near that agreement. This covers $\GL(K)$ acting on Gaussian laws away from consensus. At consensus it fails for $K\geq3$: identical Gaussian laws with nonzero mean share a stabilizer of dimension $(K-1)(K-2)/2$, so the transport Hessian is singular there.

\section{A continuum energy and its auxiliary connection}

Write $V(x,p,b)$ for the local potential $\mathcal E$ at context $x$, with rows in log-odds coordinates $b$. Supply the base $\mathcal C=\mathbb T^d=\R^d/(2\pi\mathbb Z)^d$ with its flat metric, normalized measure $d\mu=(2\pi)^{-d}d^dx$, and periodic data. Work in the trivial principal bundle, with comparisons independent of $x$ in this trivialization. For a periodic Lie-algebra-valued one-form $A=A_\mu dx^\mu$, set
\begin{equation}
 \begin{aligned}
 D^A_\mu p&=\partial_\mu p+J_pA_\mu,\\
 \mathcal F[p,b,A]&=\int_{\mathcal C}
 \left[V(x,p,b)+\frac\lambda2\sum_\mu
                  \|D^A_\mu p\|_{\mathsf G}^2\right]d\mu.
 \end{aligned}
 \label{comp:ctx:continuum-energy}
\end{equation}
where $b$ collects row log odds and $\lambda>0$. The connection has no separate kinetic term or evolution law. Assume a free proper isometric action on an invariant open region $\mathcal N_0$. Then the locked inertia $\mathfrak k_p(\xi,\eta)=\mathsf G_p(J_p\xi,J_p\eta)$ is positive, and the auxiliary equation is the strictly convex quadratic minimization
\begin{equation}
 \mathfrak k_p(A^*_\mu,\xi)
       =-\mathsf G_p(\partial_\mu p,J_p\xi),\qquad
 D^{A^*}_\mu p=P_H\partial_\mu p.
 \label{comp:ctx:continuum-normal}
\end{equation}
Here $P_H$ is Fisher-orthogonal projection off the group orbit. This is the mechanical connection, familiar from separating rotations and internal shape \citep{LittlejohnReinsch1997}. Expanding a finite alignment as $U^*(\epsilon)=\exp[-\epsilon A(X)+O(\epsilon^2)]$ gives the residual $\epsilon(dp(X)+J_pA(X))+O(\epsilon^2)$, and \cref{comp:ctx:finite-normal} reduces to \cref{comp:ctx:continuum-normal}. This local expansion establishes consistency, without supplying a graph-refinement theorem.

Use one adapted chart $p=\Psi(u,r)$, where $u$ parametrizes the quotient $\mathcal N_0/G$ and $r$ the group fiber. The metric has the form
\begin{equation}
 \mathsf G=h_{ab}(u)du^adu^b
 +k_{\alpha\beta}(u,r)
       (dr^\alpha+C^\alpha_a du^a)(dr^\beta+C^\beta_b du^b).
 \label{comp:ctx:adapted-metric}
\end{equation}
Both $h$ and $k$ are positive. The horizontal vectors are $H_a=\partial_{u^a}-C^\alpha_a\partial_{r^\alpha}$, and $h$ is the quotient metric of the Riemannian submersion \citep{ONeill1966}. Auxiliary elimination gives
\begin{equation}
 \mathcal F_*[u,r,b]=\int_{\mathcal C}
 \left[V(x,u,r,b)+\frac\lambda2h_{ab}(u)
                         \partial_\mu u^a\partial_\mu u^b\right]d\mu.
 \label{comp:ctx:reduced}
\end{equation}
Thus contextual energy sees shape variation. Pure common frame variation costs zero, and a transitive group action leaves no shape directions to regularize. This is the elimination of a nondynamical gauge field in a gauged sigma model, the mechanism of the $CP^{N-1}$ model \citep{DAddaLuscherDiVecchia1978} and of hidden local symmetry \citep{BandoKugoYamawaki1988}; the Fisher target metric, the nonbasic reference force, and the frame and row equations are what the present construction adds.

\section{Diffusion of shape coupled to frame and row motion}

For a row with $l$ entries, take $b_j=\log(\beta_j/\beta_l)$, $j<l$. Its Fisher matrix is $\mathsf B_{jk}=\beta_j\delta_{jk}-\beta_j\beta_k$; collect these positive blocks for all rows. The field metric is the integral of \cref{comp:ctx:adapted-metric} and $db^T\mathsf Bdb$. It is the Fisher metric of a normalized experiment that samples $x\sim\mu$ and then the local product laws and source labels. Define the quotient tension
\begin{equation}
 \tau_h(u)^a=\Delta u^a+
          \Gamma^a_{bc}(u)\partial_\mu u^b\partial_\mu u^c,
 \label{comp:ctx:tension}
\end{equation}
with $\Gamma$ the Levi-Civita coefficients of $h$. Periodic integration by parts gives the energy derivatives $V_u-\lambda h\tau_h$, $V_r$, and $V_b$. The metric-gradient equations are
\begin{equation}
 \begin{aligned}
 h u_t+C^Tk(r_t+Cu_t)&=-V_u+\lambda h\tau_h,\\
 k(r_t+Cu_t)&=-V_r,\qquad \mathsf B b_t=-V_b.
 \end{aligned}
\end{equation}
Solving these equations yields
\begin{equation}
 \begin{aligned}
 u_t&=\lambda\Delta u+\lambda\Gamma(u)(\partial_\mu u,\partial_\mu u)
                           -h^{-1}(V_u-C^TV_r),\\
 r_t&=-C(u,r)u_t-k(u,r)^{-1}V_r,\qquad
 b_t=-\mathsf B(b)^{-1}V_b.
 \end{aligned}
 \label{comp:ctx:coupled-pde}
\end{equation}
Fixed references may distinguish a frame direction, so $V_r$ need not vanish. The $C^TV_r$ term is then required by the Fisher metric. Transforming references, evaluators, and channels together with the laws is a passive covariance statement; it supplies no active symmetry with those data held fixed. A smooth frame change fixed in evolution time obeys $A'=\operatorname{Ad}_gA-(dg)g^{-1}$ and carries the equations covariantly. A time-dependent change adds its coordinate velocity.

The shape operator is the harmonic-map heat operator \citep{EellsSampson1964}. The complete law equation is degenerate in frame directions: its frame equation contains $u_t$ rather than an independent frame Laplacian. This is forced by gauge invariance, not by the particular energy: any contextual energy invariant under context-dependent frame changes has a horizontal Fisher gradient, so after elimination the frame velocity is the pointwise one of the local potential. For each source row of sector $a\in\{q,s\}$, \cref{comp:ctx:coupled-pde} reads
\begin{equation}
 \partial_t b_j=-(D_j-D_l)
       -T_a\left(b_j-\log\frac{\pi_j}{\pi_l}\right).
 \label{comp:ctx:log-odds}
\end{equation}
Bounded costs on a finite interval bound the log odds and keep rows positive. They do not provide a separate spatial smoothing estimate for rows. Instantaneous Gibbs rows $\beta_j^*\propto\pi_j e^{-D_j/T_a}$ instead replace their energy by $-T_a\log\sum_j\pi_j e^{-D_j/T_a}$, producing a different law-only flow.

\theoremheading{Local continuum evolution}{comp:ctx:local-theorem}
Assume the periodic geometry above, the declared regular finite-dimensional product family, a free proper action, and one adapted chart with smooth metric coefficients and finite smooth $V(x,u,r,b)$ on that open chart. Fix $\lambda>0$, an integer $s>d/2+1$, and $m=s-1$. If $u_0\in H^s(\mathbb T^d)$, $r_0,b_0\in H^m(\mathbb T^d)$, and their pointwise image lies in a compact subset of the regular interior chart, then a unique local solution satisfies
\begin{equation}
 \begin{aligned}
 u&\in C_tH^s_x\cap L^2_tH^{s+1}_x\cap H^1_tH^m_x,\\
 r,b&\in C_tH^m_x\cap H^1_tH^m_x,\qquad b\in C^1_tH^m_x.
 \end{aligned}
 \label{comp:ctx:solution-class}
\end{equation}
Its existence time depends only on $\lambda$, bounds for the data norms, and a compact set containing the image. It depends locally Lipschitz-continuously on the initial data in these norms on a common short interval, and continues while $\|u\|_{H^s}+\|r\|_{H^m}+\|b\|_{H^m}$ stays bounded and its image remains in a fixed compact subset of the chart.

\begin{proof}
Write the first equation as $u_t-\lambda\Delta u=F(z)$ with $z=(u,r,b)$, and the others as $r_t=-C(z)u_t+Q(z)$ and $b_t=L(z)$. Since $m>d/2$, products are bounded in $H^m$. Smooth composition on bounded sets with compact image is locally Lipschitz. The gradient product in $F$ also lies in $H^m$, because $u\in H^s$ gives $\partial_\mu u\in H^m$. Thus $F,C,Q,L$ are bounded and locally Lipschitz from $H^s\times H^m\times H^m$ into $H^m$.

Fourier expansion of the forced heat equation gives, for a time horizon $T\leq1$,
\begin{equation}
 \|v\|_{C_tH^s_x}+\|v\|_{L^2_tH^{s+1}_x}
       +\|v_t\|_{L^2_tH^m_x}
 \leq K_\lambda\bigl(\|v_0\|_{H^s}+\|f\|_{L^2_tH^m_x}\bigr).
 \label{comp:ctx:heat-estimate}
\end{equation}
For each nonzero Fourier mode, multiply its equation by the $H^s$ weight times its conjugate and apply Young's inequality; $|k|^2(1+|k|^2)^s$ controls the $H^{s+1}$ weight. The zero mode is bounded by integrating its forcing. The equation itself bounds $v_t$, proving the estimate. Fourier approximation gives continuity at the initial time.

On continuous paths in $H^s\times H^m\times H^m$, first solve the heat equation with forcing $F(z)$ to obtain $u^+$. Then define
\begin{equation}
 r^+(t)=r_0+\int_0^t[-C(z)u_t^++Q(z)]da,\qquad
 b^+(t)=b_0+\int_0^tL(z)da.
 \label{comp:ctx:iteration}
\end{equation}
The new heat derivative appears here, so the input path needs no time derivative. The heat estimate bounds $\|u_t^+\|_{L^2H^m}$ by a constant $A_T$ bounded as $T\downarrow0$. For two inputs at distance $D$ in the continuous-path norm, the zero-initial heat estimate bounds their output heat difference and derivative difference by $K_\lambda L_F\sqrt T D$. Splitting the fiber difference as $C(z)(u_t^+-\widetilde u_t^+)+(C(z)-C(\widetilde z))\widetilde u_t^+$ then bounds it by a constant times $(T+\sqrt T A_T)D$; the row difference is bounded by $TL_LD$. The map is therefore a contraction for small $T$.

It preserves a bounded path set with small displacement in $H^m$: the homogeneous heat displacement is at most $\sqrt{\lambda T}\|u_0\|_{H^s}$, the forced displacement is $O(T)$, and the fiber displacement is $O(\sqrt T A_T+T)$. Sobolev embedding preserves the interior chart. The fixed point has \cref{comp:ctx:solution-class} by the heat estimate and the two integral equations. Keeping initial-data differences proves the stated dependence and uniqueness. Bounded norms and a fixed compact image make all these constants uniform, giving a uniform restart interval and hence the continuation assertion.
\end{proof}

The energy chain rule is valid in \cref{comp:ctx:solution-class}. Pairing the gradient equations with their velocities gives, almost everywhere in time,
\begin{equation}
 \frac{d\mathcal F_*}{dt}
 =-\int_{\mathcal C}
 \left[u_t^Thu_t+V_r^Tk^{-1}V_r+V_b^TB^{-1}V_b\right]d\mu.
 \label{comp:ctx:continuum-dissipation}
\end{equation}
The energy is absolutely continuous and the integrated identity holds from the initial time. Its decrease controls the total energy, without supplying the higher Sobolev bounds in the continuation criterion. Every estimate uses only $m>d/2$, through the algebra property, the embedding into continuous functions, and smooth composition \citep{Taylor2011}. Setting $\lambda=0$ instead gives a locally Lipschitz ODE in $H^m$ for all coordinates. The extra derivative assigned to $u$ need not persist in that ODE, and the positive-coupling heat estimates are not uniform as $\lambda\downarrow0$.

\section{An exact amplitude that diffuses across contexts}

Take two agents on $\mathcal C=S^1$, with $q_i=\mathcal N(\mu_i,I_2)$, $s_i=\mathcal N(\nu_i,I_2)$, $r_i=\mathcal N(0,I_2)$, and $G_i(m,dk)=\mathcal N(m,I_2)(dk)$. Each agent attends only to the other through an identity channel in each sector. The rows are therefore singleton rows. Direct integration gives the local contribution $\frac12\|\mu_i-\nu_i\|^2+\frac12\|\nu_i\|^2+1$, so both sectors participate. With $P=(\mu_1,\nu_1,\mu_2,\nu_2)$ the full reciprocal potential is
\begin{equation}
 V(P)=2+\frac12\langle P,P\mathsf H\rangle_F,\qquad
 \mathsf H=\begin{pmatrix}
 3&-1&-2&0\\-1&4&0&-2\\-2&0&3&-1\\0&-2&-1&4
 \end{pmatrix}.
 \label{comp:ctx:gaussian-potential}
\end{equation}
The Fisher metric is Frobenius. A common $SO(2)$ rotation acts on all columns and is free whenever $P\ne0$. For $J=\left(\begin{smallmatrix}0&-1\\1&0\end{smallmatrix}\right)$, write $A_x=\alpha J$. The auxiliary equation and law flow become
\begin{equation}
 \alpha_*=-\frac{\langle P_x,JP\rangle_F}{\|P\|_F^2},\qquad
 P_t=-P\mathsf H+\lambda(\partial_x+\alpha_*J)^2P.
 \label{comp:ctx:matrix-flow}
\end{equation}
The derivative in the last expression acts on $\alpha_*$ as well as $P$. It follows by first varying with $\alpha$ fixed, using skew adjointness of $\partial_x+\alpha J$, and then substituting its auxiliary minimizer.

Put $\varsigma=(\sqrt5-1)/2$, $\kappa=5-\varsigma$, $v=(1,\varsigma,-1,-\varsigma)^T$, and $E_0=e_1v^T$. Let $R_\varphi=\exp(\varphi J)$. Since $\mathsf Hv=\kappa v$, the field $P(t,x)=R_{\varphi(x)}a(t,x)E_0$, with any smooth periodic phase $\varphi$ and positive amplitude $a$, obeys
\begin{equation}
 \alpha_*=-\varphi_x,\qquad
 a_t=\lambda a_{xx}-\kappa a,\qquad \varphi_t=0.
 \label{comp:ctx:amplitude}
\end{equation}
The connection removes phase variation from the spatial derivative. The remaining amplitude is gauge invariant: $a=\|P\|_F/\|E_0\|_F$. For $n\geq1$ and $a_0>|a_1|>0$, an exact solution is
\begin{equation}
 a(t,x)=a_0e^{-\kappa t}
              +a_1e^{-(\kappa+\lambda n^2)t}\cos(nx).
 \label{comp:ctx:exact-amplitude}
\end{equation}
It remains positive at every finite time. Relative to the constant mode, the oscillatory mode decays by $e^{-\lambda n^2t}$; without contextual coupling that ratio stays fixed. Initial profiles with frequencies one and two have identical laws and first derivatives at $x=0$, but their local velocities differ by $-3\lambda a_1R_{\varphi(0)}E_0$. Neighboring contextual variation has become a force. The statistical shape approaches the singular zero-mean orbit only at infinite time in this explicit solution. With a fixed reference that breaks the rotation symmetry, the same family can instead reach the zero-mean orbit in finite time with bounded norms and decreasing energy; the law flow passes through it smoothly while the quotient description breaks down.

This example uses a restricted Gaussian location family and singleton rows; it supplies no invariant Gaussian family for the full density flow and no test of nontrivial row dynamics. It is the linear radial sector of the system: the quotient metric is flat along its ray, the references are rotation invariant so the nonbasic force vanishes, the phase does not move, and the connection is pure gauge with trivial holonomy around the circle. It therefore tests none of the nonlinear, frame, or nonbasic parts of the theorem. The continuum theorem concerns a supplied torus, one regular chart, fixed comparisons, and a fixed agent census. It leaves unrestricted density families, singular orbit changes, nontrivial base bundles, boundary problems, graph refinement, and global continuation open. Independent rows and communication links acquire no separate diffusion, and pure common frame variation remains unpenalized. \Cref{comp:link-base-transport} separately constructs contextual displacement transports from regular law, interface, and communication-link data. Its comparison rule need not minimize the Fisher energy used here, so its substitution would require a new energy variation and evolution analysis.
